\documentclass[12pt,a4paper]{article}

% --- Настройка языка и шрифтов для LuaLaTeX ---
\usepackage{fontspec}
\defaultfontfeatures{Ligatures=TeX,Scale=MatchLowercase}

\setmainfont{Alice-Regular}[
    Path=font/,
    Extension=.ttf,
    BoldFont = Alice-Regular,
    ItalicFont = Alice-Regular,
    BoldItalicFont = Alice-Regular,
    BoldFeatures = {FakeBold=1.6},
    ItalicFeatures = {FakeSlant=0.2},
    BoldItalicFeatures = {FakeBold=1.6,FakeSlant=0.2}
]

\newfontfamily\cyrillicfont{Alice-Regular}[
    Path=font/,
    Extension=.ttf,
    BoldFont = Alice-Regular,
    ItalicFont = Alice-Regular,
    BoldItalicFont = Alice-Regular,
    BoldFeatures = {FakeBold=1.6},
    ItalicFeatures = {FakeSlant=0.2},
    BoldItalicFeatures = {FakeBold=1.6,FakeSlant=0.2}
]

\usepackage{polyglossia}
\setmainlanguage{russian}
\setsansfont{TeX Gyre Heros}
\setmonofont{TeX Gyre Cursor}
\AtBeginDocument{\renewcommand{\textnumero}{{\sffamily №}}}

\usepackage{graphicx}                           % Графика
\usepackage{amsmath, amsfonts, amssymb, amsthm} % Математика
\usepackage{mathtools}
\usepackage{chemfig}
\usepackage[version=4,arrows=pgf-filled,
            textfontname=sffamily,
            mathfontname=mathsf]{mhchem}
\usepackage[protrusion=false]{microtype}       % Микротипографика
\usepackage{braket}                         % Braket нотация для квантовой механики
\usepackage{unicode-math}                   % Современный пакет для математики в LuaLaTeX
\usepackage{longtable}                      % Таблицы с разрывом страниц
\usepackage{tikz}                           % Рисование
\usepackage[europeanresistors]{circuitikz}  % Рисование электрических схем в европейском стиле
\usepackage{listings}                       % Для отображения кода
\usepackage{tabularx}                       % Для таблиц во всю ширину
\usepackage{array}                          % Для расширенных возможностей таблиц
\usepackage{float}                          % Для точного позиционирования фигур
\usepackage{pgfplots}                          % Для построения графиков
\usepackage{caption}                           % Для настройки подписей к рисункам и таблицам
\usepackage{xfp}                               % Для вычислений в LaTeX
\usepackage{luacode}                           % Для выполнения Lua-кода в документе    
\pgfplotsset{compat=1.18}

\usepackage[paper=a4paper,
            top=20mm,
            bottom=20mm,
            left=10mm,
            right=10mm,
            footskip=0.4in 
]{geometry}

\usetikzlibrary{arrows.meta}

% Настройка формата подписей в соответствии с образцом (ГОСТ: «Рисунок 1 — ...»)
\DeclareCaptionLabelFormat{fullfig}{Рисунок~#2}
\DeclareCaptionLabelFormat{fulltab}{Таблица~#2}
\captionsetup[figure]{labelformat=fullfig, labelsep=endash, justification=centering}
\captionsetup[table]{labelformat=fulltab, labelsep=endash, justification=justified}

% --- Колонтитулы ---
\usepackage{fancyhdr}
\pagestyle{fancy}
\fancyhf{}
\setlength{\headheight}{15pt}

% --- Теоремы и окружения ---
\theoremstyle{plain}
\newtheorem{theorem}{Theorem}[section]
\newtheorem{lemma}[theorem]{Lemma}
\newtheorem{proposition}[theorem]{Proposition}
\newtheorem{corollary}[theorem]{Corollary}

\theoremstyle{definition}
\newtheorem{definition}[theorem]{Definition}

\theoremstyle{remark}
\newtheorem{remark}[theorem]{Remark}

\numberwithin{equation}{section}

\renewcommand{\headrulewidth}{0.4pt}
\renewcommand{\footrulewidth}{0.4pt}
\renewcommand\lstlistingname{Листинг}

\lstdefinelanguage{spice}{
    alsoletter={.},
    morecomment=[l]{*},
    morecomment=[l]{;},
    morekeywords={.DC, .dc, .TRAN, .tran, .AC, .ac, .OP, .op, .SAVE, .save, .PROBE, .probe, .END, .end, .PARAM, .param, .MODEL, .model, .SUBCKT, .subckt, .ENDS, .ends, .INCLUDE, .include, .LIB, .lib, .IC, .ic, .NODESET, .nodeset, .STEP, .step, .PRINT, .print, .PLOT, .plot},
    sensitive=false
}

\lstdefinestyle{main}{
    backgroundcolor=\color{gray!8},   
    commentstyle=\color{green!50!black}\itshape,
    keywordstyle=\color{blue!80!black}\bfseries,
    numberstyle=\small\color{darkgray},
    stringstyle=\rmfamily\color{orange!90!black},
    basicstyle=\ttfamily\small,
    breakatwhitespace=false,         
    breaklines=true,                   
    keepspaces=true,                 
    numbers=left,                    
    numbersep=8pt,                  
    showspaces=false,                
    showstringspaces=false,
    showtabs=false,                  
    tabsize=4,
    frame=single,
    framerule=0.5pt,
    rulecolor=\color{gray!30},
    captionpos=t
}   

\lstset{style=main}

\newcolumntype{C}{>{\centering\arraybackslash}X}
\renewcommand{\arraystretch}{1.25}

\begin{luacode*}
    -- data = require("data")
    -- include = require("include")
    -- handle = require("handle")
\end{luacode*}

\newcommand{\VariantN}{1}
\newcommand{\LabNumber}{2}
\newcommand{\LabTitle}{Работа с графическим редактором схем и постпроцессором.\\
Система схемотехнического проектирования Symica DE}
\newcommand{\Discipline}{Разработка САПР}

\title{
    Министерство науки и высшего образования 
    
    Российской Федерации
    
    Национальный исследовательский университет «МИЭТ»
    
    Институт интегральной электроники имени академика 
    
    К. А. Валиева

    \vspace{2em}
}
\author{}
\date{}

\begin{document}

\maketitle
\thispagestyle{fancy}

\begin{center}
\Large Отчёт по лабораторной работе \textnumero~$\LabNumber$\\
по дисциплине «\Discipline»\\
«\LabTitle»
\end{center}

\vspace{12em}

\begin{flushright}
    Выполнил студент группы ЭН-$35$\\
    Андреев И. А.\\
    Вариант \textnumero~$\VariantN$\\
    \vspace{1em}
    Проверил преподаватель института ИнЭл\\
    Щучкин Евгений Юрьевич
\end{flushright}

\cfoot{Москва 2026}

\pagebreak

\fancyhead[L]{\LabTitle}
\fancyhead[R]{\Discipline}
\fancyfoot[C]{\thepage}

\begin{center}
\Large\textbf{Лабораторная работа $\LabNumber$}
\end{center}

\textbf{Цель работы:} получить навыки работы с графическим редактором схем и
графическим постпроцессором САПР Symica DE.

\begin{center}
\LARGE\textbf{\LabTitle}
\end{center}

% =========================================================================
% 1. ПРИМЕР 1
% =========================================================================
\section{Пример 1. Ток через резистор (анализ по постоянному току)}

Схема приведена на рисунке~\ref{fig:task1_circuit}, результат 
моделирования — на рисунке~\ref{fig:task1_plot}.

\begin{figure}[H]
    \centering
    \begin{tikzpicture}
        \draw[dashed] (0,0) rectangle (10,6);
        \node at (5,3) {\color{gray}вставьте снимок схемы из редактора Symica DE};
    \end{tikzpicture}
%   \includegraphics[width=0.7\textwidth]{task1_circuit.png}
    \caption{Схема из источника напряжения $V_0$ и резистора $R_0$}
    \label{fig:task1_circuit}
\end{figure}

\begin{figure}[H]
    \centering
    \begin{tikzpicture}
        \draw[dashed] (0,0) rectangle (10,6);
        \node at (5,3) {\color{gray}вставьте снимок графика из SymProbe};
    \end{tikzpicture}
%   \includegraphics[width=0.7\textwidth]{task1_plot.png}
    \caption{Ток через резистор $R_0$ при изменении напряжения источника от $0$ до $5$ В}
    \label{fig:task1_plot}
\end{figure}

\textbf{Результат:} график — прямая из начала координат с $1$ мА при $5$ В (закон Ома подтверждён).

% =========================================================================
% 2. ПРИМЕР 2
% =========================================================================
\section{Пример 2. КМОП-инвертор (анализ переходных процессов), измерение задержек}

Схема инвертора приведена на рисунке~\ref{fig:task2_circuit}, 
временные диаграммы при нагрузочных ёмкостях $10$ пФ и $100$ пФ — 
на рисунках~\ref{fig:task2_plot_10pF} и~\ref{fig:task2_plot_100pF}, 
измеренные задержки — в таблице~\ref{tab:task2_delays}.

\begin{figure}[H]
    \centering
    \begin{tikzpicture}
        \draw[dashed] (0,0) rectangle (10,6);
        \node at (5,3) {\color{gray}вставьте снимок схемы из редактора Symica DE};
    \end{tikzpicture}
%   \includegraphics[width=0.7\textwidth]{task2_circuit.png}
    \caption{КМОП-инвертор с ёмкостной нагрузкой}
    \label{fig:task2_circuit}
\end{figure}

\begin{figure}[H]
    \centering
    \begin{tikzpicture}
        \draw[dashed] (0,0) rectangle (10,6);
        \node at (5,3) {\color{gray}вставьте снимок графика из SymProbe};
    \end{tikzpicture}
%   \includegraphics[width=0.7\textwidth]{task2_plot_10pF.png}
    \caption{Временная диаграмма работы инвертора при ёмкости $10$ пФ}
    \label{fig:task2_plot_10pF}
\end{figure}

\begin{figure}[H]
    \centering
    \begin{tikzpicture}
        \draw[dashed] (0,0) rectangle (10,6);
        \node at (5,3) {\color{gray}вставьте снимок графика из SymProbe};
    \end{tikzpicture}
%   \includegraphics[width=0.7\textwidth]{task2_plot_100pF.png}
    \caption{Временная диаграмма работы инвертора при ёмкости $100$ пФ}
    \label{fig:task2_plot_100pF}
\end{figure}

\begin{table}[H]
    \centering
    \begin{tabularx}{\textwidth}{|C|C|C|}
        \hline
        Нагрузочная ёмкость & Задержка $t_{PLH}$, нс & Задержка $t_{PHL}$, нс \\
        \hline
        $10$ пФ & & \\ \hline
        $100$ пФ & & \\
        \hline
    \end{tabularx}
    \caption{Измеренные задержки переключения инвертора}
    \label{tab:task2_delays}
\end{table}

\textbf{Результат:} при увеличении ёмкости с $10$ до $100$ пФ задержка выросла примерно в~$\dots$ раз(а).

% =========================================================================
% 3. ПРИМЕР 3
% =========================================================================
\section{Пример 3. Элемент $2$-И-НЕ (NAND$2$)}

Схема элемента приведена на рисунке~\ref{fig:task3_circuit}, временная диаграмма — 
на рисунке~\ref{fig:task3_plot}, снятая по ней таблица истинности — в таблице~\ref{tab:task3_truth_table}.

\begin{figure}[H]
    \centering
    \begin{tikzpicture}
        \draw[dashed] (0,0) rectangle (10,6);
        \node at (5,3) {\color{gray}вставьте снимок схемы из редактора Symica DE};
    \end{tikzpicture}
%   \includegraphics[width=0.7\textwidth]{task3_circuit.png}
    \caption{Схема электрическая принципиальная элемента $2$-И-НЕ (NAND$2$)}
    \label{fig:task3_circuit}
\end{figure}

\begin{figure}[H]
    \centering
    \begin{tikzpicture}
        \draw[dashed] (0,0) rectangle (10,6);
        \node at (5,3) {\color{gray}вставьте снимок графика из SymProbe};
    \end{tikzpicture}
%   \includegraphics[width=0.7\textwidth]{task3_plot.png}
    \caption{Временная диаграмма работы элемента $2$-И-НЕ (NAND$2$)}
    \label{fig:task3_plot}
\end{figure}

\begin{table}[H]
    \centering
    \begin{tabularx}{\textwidth}{|C|C|C|}
        \hline
        Вход А & Вход В & Выход out по диаграмме \\
        \hline
        0 & 0 & \\ \hline
        0 & 1 & \\ \hline
        1 & 0 & \\ \hline
        1 & 1 & \\
        \hline
    \end{tabularx}
    \caption{Таблица истинности элемента $2$-И-НЕ по результатам моделирования}
    \label{tab:task3_truth_table}
\end{table}

\textbf{Результат:} снятая диаграмма совпала с таблицей истинности 
$2$-И-НЕ (NAND$2$): выход $0$ формируется только при комбинации 
$A=1$ и $B=1$.

% =========================================================================
% 4. ЗАДАНИЕ 1
% =========================================================================
\section{Задание 1. Резистивный делитель (анализ по постоянному току)}

Параметры варианта и результат моделирования приведены в таблице~\ref{tab:task4_variant}, 
схема — на рисунке~\ref{fig:task4_circuit}, графики — на рисунке~\ref{fig:task4_plot}.

\begin{table}[H]
    \centering
    \begin{tabularx}{\textwidth}{|C|C|C|C|C|}
        \hline
        $R_1$, кОм & $R_2$, кОм & Развёртка источника (START\dots STOP, STEP) & $V_{\text{out}}$ при $V_1 = \text{STOP}$, В & Ток в цепи, мА \\
        \hline
        $1$ & $1$ & $0\dots 5$~В, $10$~мВ & & \\
        \hline
    \end{tabularx}
    \caption{Вариант задания $1$ и результат моделирования}
    \label{tab:task4_variant}
\end{table}

\begin{figure}[H]
    \centering
    \begin{tikzpicture}
        \draw[dashed] (0,0) rectangle (10,6);
        \node at (5,3) {\color{gray}вставьте снимок схемы из редактора Symica DE};
    \end{tikzpicture}
%   \includegraphics[width=0.7\textwidth]{task4_circuit.png}
    \caption{Схема резистивного делителя напряжения}
    \label{fig:task4_circuit}
\end{figure}

\begin{figure}[H]
    \centering
    \begin{tikzpicture}
        \draw[dashed] (0,0) rectangle (10,6);
        \node at (5,3) {\color{gray}вставьте снимок графика из SymProbe};
    \end{tikzpicture}
%   \includegraphics[width=0.7\textwidth]{task4_plot.png}
    \caption{Ток и напряжения делителя при изменении напряжения источника}
    \label{fig:task4_plot}
\end{figure}

\textbf{Результат:} полученные значения $V_{\text{out}}$ и тока в цепи совпали со значениями по варианту.

% =========================================================================
% 5. ЗАДАНИЕ 2
% =========================================================================
\section{Задание 2. Подбор резистора по задержке (анализ переходных процессов)}

Параметры варианта, подобранное сопротивление и измеренная задержка 
приведены в таблице~\ref{tab:task5_variant}, схема — на рисунке~\ref{fig:task5_circuit}, 
временная диаграмма с курсорами — на рисунке~\ref{fig:task5_plot}.

\begin{table}[H]
    \centering
    \begin{tabularx}{\textwidth}{|C|C|C|C|C|}
        \hline
        Ёмкость $C_1$ по варианту & Целевая задержка, пс & Подобранное сопротивление $R$ & Измеренная задержка, пс & Отклонение, пс \\
        \hline
        $1$ пФ & $400$ пс & & & \\
        \hline
    \end{tabularx}
    \caption{Вариант задания 2 и результат подбора}
    \label{tab:task5_variant}
\end{table}

\begin{figure}[H]
    \centering
    \begin{tikzpicture}
        \draw[dashed] (0,0) rectangle (10,6);
        \node at (5,3) {\color{gray}вставьте снимок схемы из редактора Symica DE};
    \end{tikzpicture}
%   \includegraphics[width=0.7\textwidth]{task5_circuit.png}
    \caption{Схема интегрирующей RC-цепи}
    \label{fig:task5_circuit}
\end{figure}

\begin{figure}[H]
    \centering
    \begin{tikzpicture}
        \draw[dashed] (0,0) rectangle (10,6);
        \node at (5,3) {\color{gray}вставьте снимок графика из SymProbe};
    \end{tikzpicture}
%   \includegraphics[width=0.7\textwidth]{task5_plot.png}
    \caption{Измерение задержки переключения по уровню $50\%$}
    \label{fig:task5_plot}
\end{figure}

\textbf{Результат:} измеренная задержка уложилась в допустимый диапазон $\pm 10$~пс.

% =========================================================================
% 6. ЗАДАНИЕ 3
% =========================================================================
\section{Задание 3. Элемент $2$-ИЛИ-НЕ (NOR$2$)}

Схема элемента приведена на рисунке~\ref{fig:task6_circuit}, временная диаграмма — 
на рисунке~\ref{fig:task6_plot}, снятая по ней таблица истинности — в таблице~\ref{tab:task6_truth_table}.

\begin{figure}[H]
    \centering
    \begin{tikzpicture}
        \draw[dashed] (0,0) rectangle (10,6);
        \node at (5,3) {\color{gray}вставьте снимок схемы из редактора Symica DE};
    \end{tikzpicture}
%   \includegraphics[width=0.7\textwidth]{task6_circuit.png}
    \caption{Схема электрическая принципиальная элемента $2$-ИЛИ-НЕ (NOR$2$)}
    \label{fig:task6_circuit}
\end{figure}

\begin{figure}[H]
    \centering
    \begin{tikzpicture}
        \draw[dashed] (0,0) rectangle (10,6);
        \node at (5,3) {\color{gray}вставьте снимок графика из SymProbe};
    \end{tikzpicture}
%   \includegraphics[width=0.7\textwidth]{task6_plot.png}
    \caption{Временная диаграмма работы элемента $2$-ИЛИ-НЕ (NOR$2$)}
    \label{fig:task6_plot}
\end{figure}

\begin{table}[H]
    \centering
    \begin{tabularx}{\textwidth}{|C|C|C|}
        \hline
        Вход А & Вход В & Выход out по диаграмме \\
        \hline
        0 & 0 & \\ \hline
        0 & 1 & \\ \hline
        1 & 0 & \\ \hline
        1 & 1 & \\
        \hline
    \end{tabularx}
    \caption{Таблица истинности элемента $2$-ИЛИ-НЕ по результатам моделирования}
    \label{tab:task6_truth_table}
\end{table}

\textbf{Результат:} снятая диаграмма полностью совпала с таблицей 
истинности логического элемента $2$-ИЛИ-НЕ (NOR$2$): выход равен $1$ 
только при $A=0$ и $B=0$.

% =========================================================================
% 7. ВЫВОДЫ ПО РАБОТЕ
% =========================================================================
\section{Выводы по работе}

В ходе выполнения лабораторной работы были освоены базовые приёмы 
работы в графическом схемотехническом редакторе Symica DE и 
постпроцессоре SymProbe. Были проведены статический расчёт режимов по 
постоянному току (DC) и анализ динамических характеристик во временной 
области (TRAN) для пассивных цепей и КМОП-логических вентилей. 
Экспериментально подтверждена пропорциональная зависимость задержки 
переключения КМОП-инвертора от величины нагрузочной ёмкости, а также 
верифицирована логика функционирования базовых элементов $2$-И-НЕ и $2$-ИЛИ-НЕ.

\end{document}